% !TEX root = ../main.tex
\chapter{Wstęp}
\label{ch:wstep}

% Zgodnie z wytycznymi szablonu wstęp powinien zawierać:
% - wprowadzenie w problem/zagadnienie
% - osadzenie problemu w dziedzinie
% - cel pracy
% - zakres pracy
% - zwięzłą charakterystykę rozdziałów

\section{Wprowadzenie do tematu}

Platformy marketplace to internetowe serwisy pośredniczące w~sprzedaży,
w~których operator zapewnia infrastrukturę techniczną, komunikację i~mechanizmy
bezpieczeństwa, natomiast transakcje są zawierane między użytkownikami.
W~modelu C2C, czyli konsument do konsumenta, sprzedawcą i~kupującym są osoby
prywatne. Takie platformy przeżywają dynamiczny rozwój na~całym świecie.
W~Polsce serwisy takie jak Allegro, OLX czy~Vinted należą do~najszerzej
używanych kanałów handlu między konsumentami i~generują wolumen transakcji,
który czyni je~atrakcyjnym celem dla potencjalnych oszukujących \cite{bib:mayayise2024}.

Nadużycia na~platformach marketplace stanowią poważne wyzwanie zarówno dla
operatorów, jak i~dla użytkowników. Obejście platformy, wyłudzenia oraz toksyczność
w~komunikacji mogą prowadzić do~utraty zaufania, znaczących spadków przychodów
oraz negatywnych doświadczeń użytkowników. W~dobie rosnącej popularności takich
platform skuteczne wykrywanie i~zapobieganie nadużyciom staje się kluczowym
elementem strategii zarządzania ryzykiem oraz finansami operatora.

Skala problemu jest spora. Lin, Nian i~Foutz \cite{bib:lin2022airbnb} ustalili,
że~na~Airbnb 5,4\% transakcji w~Austin w~2019~roku było realizowanych poza platformą,
a~na~chińskiej platformie ZBJ przychody usługodawców wykraczały aż dziesięciokrotnie
ponad wartości zarejestrowane w~systemie. Operatorzy próbują temu zapobiegać
poprzez ukrywanie danych kontaktowych w~interfejsie, jednak jak pokazują te~badania,
skutecznym środkiem zaradczym okazuje się dopiero aktywna analiza treści wiadomości.

Automatyczne systemy moderacji są niezbędne, ponieważ ręczna weryfikacja każdej
wiadomości jest niewykonalna przy setkach tysięcy dziennych interakcji. Konieczność
działania w~czasie rzeczywistym, niska tolerancja na~fałszywe alarmy blokujące
legalne transakcje oraz specyfika kolokwialnego języka czatu czynią z~tego
problemu interesujące wyzwanie inżynieryjne i~badawcze.

\section{Osadzenie problemu w dziedzinie}

Problem detekcji nadużyć w~tekstach czatowych należy do szerszego obszaru
wykrywania nadużyć i oszustw w~systemach internetowych. Jest jednocześnie
zagadnieniem w~ramach przetwarzania języka naturalnego (Natural Language
Processing, NLP) oraz bezpieczeństwa i~zarządzania ryzykiem.

Dziedzina zarządzania bezpieczeństwem i~zaufaniem (Trust \& Safety)
wykształciła się w~ostatnich dwóch dekadach jako dyscyplina łącząca metody
uczenia maszynowego, inżynierii oprogramowania i~polityki moderacji. Pierwotnie
skupiała się na~moderacji treści w~mediach społecznościowych, stopniowo rozszerzając
zasięg na~platformy transakcyjne, gdzie do~problemu toksyczności dołączają
wektory ataku specyficzne dla~handlu: obejście prowizji, wyłudzenia finansowe
i~phishing, czyli podszywanie się pod zaufane instytucje w~celu wyłudzenia danych
lub środków finansowych \cite{bib:phishing}.

W~ostatnich latach zaobserwowano znaczący postęp w~dziedzinie NLP dzięki modelom
transformerowym - zaawansowana architektura sieci neuronowej stworzona do analizowania sekwencji danych (np. słów w zdaniu), która zamiast czytać tekst słowo po słowie, odczytuje całe wypowiedzi naraz, używając mechanizmu uwagi do wyłapywania tego, które słowa są ze sobą powiązane i najważniejsze. Model BERT \cite{bib:bert}, a~następnie jego warianty wielojęzyczne
takie jak XLM-RoBERTa \cite{bib:xlmr}, zrewolucjonizowały podejście do wielu zadań
klasyfikacji tekstu \cite{bib:transformers_survey}. W~obszarze moderacji treści
transformery osiągają wyniki niedostępne dla wcześniejszych podejść opartych
wyłącznie na~regułach czy~metodach statystycznych, szczególnie
w~zadaniach wymagających rozumienia kontekstu semantycznego i~intencji nadawcy.

Jednak większość badań nad klasyfikacją toksyczności lub detekcją
oszustw prowadzono dla języka angielskiego. Polska literatura naukowa na ten temat
pozostaje ograniczona, a~modele dostosowane do specyfiki polskiego tekstu czatowego
(kolokwializmy, skróty, błędy ortograficzne) są rzadkie.

Stąd wynika uzasadnienie dla pracy łączącej dostrojenie modelu polskojęzycznego,
porównanie z~podejściami hybrydowymi i~regułowymi oraz ewaluację na~realnych
danych z~platformy marketplace, które będą napływały z~aplikacji Stylify.

\section{Cel i~zakres pracy}
\label{sec:cel}

Celem pracy jest implementacja, porównanie i~ocena czterech podejść do automatycznej
detekcji nadużyć w~polskojęzycznym czacie platformy C2C na przykładzie aplikacji
Stylify, aplikacji zaimplementowanej przez autora tej pracy w ramach pracy inżynierskiej. Realizacja celu obejmuje:

\begin{enumerate}
  \item zaprojektowanie i~wdrożenie czterech modułów detekcji: systemu regułowego
    (baseline), hybrydowego wykrywania danych osobowych z~użyciem systemu Presidio
    i~wyrażeń regularnych, klasyfikatora opartego na dostrojonym modelu HerBERT,
    oraz kombinacji hybrydowej Presidio~+~HerBERT;
  \item przygotowanie zbioru danych na podstawie realnego ruchu komunikacyjnego
    platformy, wraz z panelem moderacji do ręcznej adnotacji wiadomości, których nadawcy zostali zidentyfikowani jako potencjalni nadużywający;
  \item ewaluację porównawczą wszystkich podejść pod kątem precyzji, pełności,
    miary~$F_1$ oraz odsetka fałszywych alarmów;
  \item integrację najskuteczniejszego rozwiązania z~aplikacją produkcyjną
    wraz z~mechanizmem tury konwersacyjnej wykrywającym ataki fragmentowane, o których mowa ~~~w podrozdziale \ref {sec:tura}.
\end{enumerate}

Zakres pracy ogranicza się do polskojęzycznej komunikacji tekstowej w~czacie C2C.
Analizowane są trzy kategorie nadużyć: obejście systemu prowizyjnego platformy,
oszustwa finansowe oraz toksyczność. Praca nie obejmuje analizy obrazów, dźwięku
ani komunikacji wielojęzycznej.

\section{Pytanie badawcze}
\label{sec:pytanie}

Główne pytanie badawcze niniejszej pracy brzmi:

\begin{quote}
\textit{Które z~podejść: regułowe oparte na detekcji danych osobowych, neuronowe z~douczaniem
polskiego modelu językowego czy hybrydowe łączące detekcję danych osobowych z~klasyfikatorem ML,
najskuteczniej wykrywa nadużycia w~polskojęzycznym czacie platformy marketplace
przy jednoczesnym zachowaniu akceptowalnej liczby fałszywych alarmów?}
\end{quote}

\section{Wkład autora}
\label{sec:wklad}

Autor niniejszej pracy samodzielnie:

\begin{itemize}
  \item zaprojektował i~zaimplementował architekturę serwisu ML jako osobnego
    mikroserwisu (FastAPI), zintegrowanego z~backendowym serwerem
    aplikacji Stylify;
  \item opracował cztery moduły detekcji nadużyć (A~-~D) opisane w~rozdziale~3.3,
    w~tym regułowe wzorce obejścia platformy oraz mechanizm hybrydowego
    potoku klasyfikacji;
  \item utworzył sztuczny zbiór danych oraz panel moderacji do ręcznej adnotacji wiadomości, których nadawcy zostali zidentyfikowani jako potencjalni oszukujący. Panel ten umożliwia moderatorom przeglądanie wiadomości, ich kontekstu (poprzednie wiadomości tej samej tury konwersacyjnej) oraz przypisywanie etykiet nadużycia;
  \item przeprowadził doszkalanie modelu HerBERT \cite{bib:herbert}
    na~zbiorze własnym obejmującym cztery klasy: \textit{bypass}, \textit{fraud},
    \textit{toxic} i~\textit{benign}, opisane w~rozdziale~3.2;
  \item zaprojektował mechanizm tury konwersacyjnej łączący kolejne wiadomości
    tego samego nadawcy w~jeden kontekst klasyfikacji, co umożliwia wykrywanie
    ataków rozbitych na wiele krótkich wiadomości;
  \item przeprowadził ewaluację porównawczą wszystkich podejść, opracował potok uczenia maszynowego (ang. pipeline) najskuteczniejszego rozwiązania do integracji z~aplikacją produkcyjną.
\end{itemize}

\section{Charakterystyka rozdziałów}
\label{sec:charakterystyka}

\textbf{Rozdział~1: Wstęp} wprowadza problem detekcji nadużyć na~platformach
konsument do konsumenta, osadza go w~dziedzinie rozumienia języka naturalnego oraz Zaufania \& Bezpieczeństwa (ang. Trust \& Safety), formułuje cel,
zakres i~pytanie badawcze pracy oraz opisuje wkład autora.

\textbf{Rozdział~2: Analiza tematu} przedstawia stan wiedzy w~zakresie
moderacji treści internetowych. Omówione są architektura transformerów
i~polskojęzyczne modele językowe,
ze~szczególnym uwzględnieniem modelu HerBERT \cite{bib:herbert}. Rozdział zawiera
przegląd istniejących rozwiązań: systemów regułowych, komercyjnych
narzędzi do~detekcji danych osobowych i~toksyczności \cite{bib:perspective},
komercyjnych systemów zapobiegającym oszustwom oraz lukę badawczą dla~polskojęzycznych
platform C2C. Osobna sekcja poświęcona jest analizie ograniczeń podejścia
opartego na~zewnętrznych dużych modelach językowych (GPT-4, Claude, Gemini)
jako moderatorów treści, obejmującej aspekty RODO, kosztów skalowalności,
niedeterministyczności oraz podatności na~wstrzyknięcie instrukcji.

\textbf{Rozdział~3: Przedmiot pracy} opisuje architekturę systemu
i~uzasadnienie wyborów projektowych: strukturę mikroserwisu ML,
cztery porównywane moduły detekcji, mechanizm tury konwersacyjnej
oraz integrację z~aplikacją Stylify. Zawiera także pełną charakterystykę
syntetycznego zbioru danych: schemat klas, metodykę generowania, statystyki
opisowe i~podział na~podzbiory, a~także opis panelu moderatora jako mechanizmu
sukcesywnego wzbogacania zbioru o~dane produkcyjne.

\textbf{Rozdział~4: Badania} prezentuje wyniki ewaluacji porównawczej
na~120 przykładowym zbiorze testowym. Każdy moduł omawiany jest osobno,
z~analizą per-klasową i~kategoryzacją błędów. Rozdział zawiera zbiorcze wyniki, wykresy progresji skuteczności, macierze pomyłek i~analizę błędów,
a~kończy się bezpośrednią odpowiedzią na~pytanie badawcze.

\textbf{Rozdział~5: Podsumowanie} przedstawia wnioski z~badań, omawia
ograniczenia pracy (syntetyczność zbioru, brak GPU, statyczność wzorców)
i~wskazuje kierunki dalszych badań.

